\paragrafos{Κλασικός ορισμός της πιθανότητας}{Γ.3.3, «Κλασικός ορισμός της πιθανότητας»}

\begin{anagke}
Για τη θεατρική παράσταση της τάξης η κυρία Αλεξάνδρα βάζει σε μια κάλπη τα ονόματα και των $25$ παιδιών και τραβά ένα χαρτάκι για τον ρόλο του αφηγητή. «Εγώ έχω πενήντα–πενήντα», λέει ο Στέλιος, «ή θα με βγάλει ή δεν θα με βγάλει!». Η Ρήνη αντιδρά: «Αν δεν βγάζει νόημα στην πραγματικότητα δεν το δέχομαι. Με το νόμισμα έχεις πενήντα–πενήντα. Εδώ έχεις ένα χαρτάκι μέσα σε είκοσι πέντε».

Έχουν δίκιο και οι δύο ότι «ή θα με βγάλει ή όχι», αλλά αυτό δεν σημαίνει ότι οι δύο περιπτώσεις είναι το ίδιο εύκολες. Θέλουμε έναν αριθμό που να μετρά πόσο εύκολο είναι να συμβεί κάτι, ώστε να τα συγκρίνουμε.

Σε αυτή την παράγραφο μαθαίνουμε τι είναι ενδεχόμενο, πότε τα αποτελέσματα είναι ισοπίθανα και πώς υπολογίζουμε την πιθανότητα ενός ενδεχομένου με τον κλασικό ορισμό. Δείχνουμε και πώς η πιθανότητα συνδέεται με τα αποτελέσματα που βλέπουμε όταν επαναλαμβάνουμε πολλές φορές το πείραμα.
\end{anagke}

\begin{lexi}{ενδεχόμενο, ισοπίθανα, πιθανότητα, βέβαιο, αδύνατο}{ονόματα και σημασίες}
Το \lexiword{ενδεχόμενο} είναι κάτι που «ενδέχεται» να συμβεί. Η \lexiword{πιθανότητα} έρχεται από το «πιθανός». Τα \lexiword{ισοπίθανα} αποτελέσματα έχουν «ίσες πιθανότητες». Βέβαιο είναι αυτό που σίγουρα συμβαίνει και \lexiword{αδύνατο} αυτό που δεν μπορεί να συμβεί.
\end{lexi}

\begin{metafrasi}
«Θα φέρω ζυγό αριθμό». & ενδεχόμενο $\mathrm{Α}=\{2,4,6\}$\\
Όλες οι έδρες του ζαριού έχουν τις ίδιες ευκαιρίες. & ισοπίθανα αποτελέσματα\\
Πόσο πιθανό είναι να φέρω ζυγό αριθμό; & $\mathrm{Ρ}(\mathrm{Α})=\dfrac{3}{6}=\dfrac{1}{2}$\\
Συμβαίνει πάντα. & βέβαιο ενδεχόμενο, πιθανότητα $1$\\
Δεν μπορεί να συμβεί. & αδύνατο ενδεχόμενο, πιθανότητα $0$\\
\end{metafrasi}

\tmima{Η θεωρία}

\begin{orismos}{ενδεχόμενο}
\textbf{Ενδεχόμενο} ενός πειράματος τύχης είναι ένα σύνολο που περιέχει ένα ή περισσότερα στοιχεία του δειγματικού χώρου. Λέμε ότι ένα ενδεχόμενο \textbf{πραγματοποιείται} όταν το αποτέλεσμα του πειράματος είναι ένα από τα στοιχεία του.
\end{orismos}

Παράδειγμα: στη ρίψη ζαριού ο δειγματικός χώρος είναι $\{1,2,3,4,5,6\}$. Ενδεχόμενα είναι τα $\{6\}$ («φέρνω $6$»), $\{1,2\}$ («φέρνω $1$ ή $2$»), $\{2,4,6\}$ («φέρνω άρτιο») και $\{4,5,6\}$ («φέρνω πάνω από $3$»). Αν το ζάρι φέρει $4$, πραγματοποιούνται τα ενδεχόμενα «άρτιος» και «πάνω από $3$» και δεν πραγματοποιούνται τα άλλα δύο.

\begin{orismos}{ισοπίθανα αποτελέσματα}
Τα αποτελέσματα ενός δειγματικού χώρου λέγονται \textbf{ισοπίθανα} όταν όλα έχουν την ίδια πιθανότητα να εμφανιστούν. Σε ένα τίμιο (αμερόληπτο) ζάρι, ένα τίμιο νόμισμα ή μια καλά ανακατεμένη κάλπη, τα αποτελέσματα είναι ισοπίθανα.
\end{orismos}

\begin{kanonas}{κλασικός ορισμός της πιθανότητας}
Όταν τα αποτελέσματα ενός πειράματος είναι ισοπίθανα, η πιθανότητα ενός ενδεχομένου $\mathrm{Α}$ είναι
$$\mathrm{Ρ}(\mathrm{Α})=\frac{\text{πλήθος ευνοϊκών αποτελεσμάτων}}{\text{πλήθος όλων των δυνατών αποτελεσμάτων}}.$$
Ευνοϊκά αποτελέσματα είναι αυτά που ανήκουν στο ενδεχόμενο (όσα μας «βολεύουν»).
\end{kanonas}

Την πιθανότητα την γράφουμε ως κλάσμα, ως δεκαδικό ή ως ποσοστό. Για παράδειγμα, στη ρίψη νομίσματος η πιθανότητα του $\mathrm{Γ}$ είναι $\dfrac{1}{2}=0{,}5=50\%$. Στη ρίψη ζαριού η πιθανότητα να φέρω $6$ είναι $\dfrac{1}{6}=0{,}1666\ldots\approx16{,}7\%$.

\begin{idiotita}{τιμές της πιθανότητας}
Επειδή τα ευνοϊκά αποτελέσματα δεν είναι περισσότερα από τα δυνατά και δεν είναι λιγότερα από $0$, ισχύει
$$0\le\mathrm{Ρ}(\mathrm{Α})\le1.$$
\textbf{Βέβαιο ενδεχόμενο} είναι αυτό που πραγματοποιείται σε κάθε εκτέλεση του πειράματος και έχει πιθανότητα $1$ (ή $100\%$). \textbf{Αδύνατο ενδεχόμενο} είναι αυτό που δεν πραγματοποιείται σε καμία εκτέλεση και έχει πιθανότητα $0$.
\end{idiotita}

\begin{center}
\begin{tikzpicture}[x=9cm]
\draw[thick] (0,0)--(1,0);
\foreach \x/\l in {0/0,0.25/{$\frac14$},0.5/{$\frac12$},0.75/{$\frac34$},1/1}{\draw[thick] (\x,0.08)--(\x,-0.08) node[below,font=\small]{\l};}
\foreach \x/\l in {0/{0\%},0.25/{25\%},0.5/{50\%},0.75/{75\%},1/{100\%}}{\node[font=\scriptsize,gray] at (\x,-0.75) {\l};}
\node[font=\small,gwarmD,align=center] at (0,0.85) {αδύνατο\\{\scriptsize (φέρνω $7$)}};
\node[font=\small,gwarmD,align=center] at (0.5,0.85) {όσο το νόμισμα\\{\scriptsize (φέρνω Γ)}};
\node[font=\small,gwarmD,align=center] at (1,0.85) {βέβαιο\\{\scriptsize (φέρνω $<7$)}};
\foreach \x in {0,0.5,1}{\fill[gwarmD] (\x,0) circle (2.5pt);}
\end{tikzpicture}
\end{center}

Αφού τα ευνοϊκά και τα μη ευνοϊκά αποτελέσματα μαζί είναι όλα τα δυνατά, η πιθανότητα να \emph{μη} συμβεί ένα ενδεχόμενο είναι $1$ μείον την πιθανότητα να συμβεί. Αν η πιθανότητα να βρέξει είναι $0{,}3$, η πιθανότητα να μη βρέξει είναι $1-0{,}3=0{,}7$.

Ο τροχός του παρακάτω σχήματος έχει $8$ ίσα τμήματα: $3$ κόκκινα, $4$ μπλε και $1$ κίτρινο. Ο δειγματικός χώρος των χρωμάτων $\{\text{κόκκινο},\text{μπλε},\text{κίτρινο}\}$ \emph{δεν} έχει ισοπίθανα στοιχεία. Ισοπίθανα είναι τα $8$ \emph{τμήματα}. Γι' αυτό μετράμε τμήματα: $\mathrm{Ρ}(\text{κόκκινο})=\dfrac{3}{8}$.

\begin{center}
\binwheel{8}{red!55,red!55,red!55,blue!40,blue!40,blue!40,blue!40,yellow!80}
\end{center}

\begin{kanonas}{σχετική συχνότητα και πιθανότητα}
Αν επαναλάβουμε ένα πείραμα $\nu$ φορές και ένα ενδεχόμενο πραγματοποιηθεί $\kappa$ φορές, τότε η \textbf{σχετική συχνότητα} του ενδεχομένου είναι $\dfrac{\kappa}{\nu}$. Όσο περισσότερες φορές επαναλαμβάνουμε το πείραμα, τόσο η σχετική συχνότητα πλησιάζει την πιθανότητα του ενδεχομένου.
\end{kanonas}

Η τάξη έριξε ένα νόμισμα πολλές φορές και κατέγραψε πόσες φορές ήρθε Γράμματα:

\begin{center}
\renewcommand{\arraystretch}{1.25}
\begin{tabular}{lccccc}
\toprule
Ρίψεις $\nu$ & $10$ & $50$ & $100$ & $500$ & $1000$\\
Γράμματα $\kappa$ & $7$ & $28$ & $54$ & $247$ & $503$\\
Σχετική συχνότητα $\dfrac{\kappa}{\nu}$ & $0{,}7$ & $0{,}56$ & $0{,}54$ & $0{,}494$ & $0{,}503$\\
\bottomrule
\end{tabular}
\end{center}

Στις λίγες ρίψεις η σχετική συχνότητα απέχει αρκετά από το $\frac{1}{2}$. Στις πολλές ρίψεις πλησιάζει το $0{,}5$, που είναι η πιθανότητα του $\mathrm{Γ}$.

\begin{prosoxi}
\begin{itemize}[leftmargin=1.2em,itemsep=0.3ex,topsep=2pt]
\item Ο κλασικός ορισμός ισχύει μόνο όταν τα αποτελέσματα είναι \emph{ισοπίθανα}. Τα αθροίσματα $2,3,\dots,12$ δύο ζαριών δεν είναι ισοπίθανα. Ισοπίθανα είναι τα $36$ ζευγάρια.
\item Δεν αρκεί να μετράς τα \emph{διαφορετικά} αποτελέσματα. Στον τροχό του σχήματος τα χρώματα είναι $3$, αλλά η πιθανότητα του καθενός δεν είναι $\frac{1}{3}$.
\item Αν $\mathrm{Ρ}(\mathrm{Α})=\frac{1}{2}$ δεν σημαίνει ότι σε $10$ ρίψεις θα συμβεί ακριβώς $5$ φορές. Μόνο σε πολλές επαναλήψεις η σχετική συχνότητα πλησιάζει το $\frac{1}{2}$.
\item Η πιθανότητα δεν μπορεί να είναι αρνητική ή μεγαλύτερη από $1$. Αν έβγαλες $\frac{7}{5}$ ή $-0{,}2$, έκανες λάθος.
\end{itemize}
\end{prosoxi}

\tmima{Λυμένα παραδείγματα}

\paradeigma{Ένα ζάρι}
\tixero Ρίχνουμε ένα τίμιο ζάρι.\par
\tithelo Την πιθανότητα των ενδεχομένων: α) περιττός αριθμός, β) αριθμός μεγαλύτερος από $4$, γ) ο αριθμός $7$, δ) αριθμός μικρότερος από $7$.\par
\possindeo Ο δειγματικός χώρος $\{1,2,3,4,5,6\}$ έχει $6$ ισοπίθανα στοιχεία. Μετρώ τα ευνοϊκά.\par
\apantisi α) $\{1,3,5\}$: $\dfrac{3}{6}=\dfrac{1}{2}=0{,}5=50\%$. β) $\{5,6\}$: $\dfrac{2}{6}=\dfrac{1}{3}\approx0{,}333\approx33{,}3\%$. γ) Κανένα ευνοϊκό: $\dfrac{0}{6}=0$ (αδύνατο). δ) Και τα $6$ ευνοϊκά: $\dfrac{6}{6}=1=100\%$ (βέβαιο).\par
\elegxos Όλες οι πιθανότητες είναι ανάμεσα στο $0$ και στο $1$.

\paradeigma{Κλήρωση στη γιορτή}
\tixero Στη γιορτή του σχολείου πουλήθηκαν $40$ λαχνοί και κερδίζουν οι $8$.\par
\tithelo Την πιθανότητα να κερδίσει κάποιος με έναν λαχνό και να μην κερδίσει.\par
\possindeo Όλοι οι λαχνοί έχουν την ίδια πιθανότητα να κληρωθούν.\par
\apantisi $\mathrm{Ρ}(\text{κερδίζω})=\dfrac{8}{40}=\dfrac{1}{5}=0{,}2=20\%$. Μη κερδίσω: $\dfrac{40-8}{40}=\dfrac{32}{40}=\dfrac{4}{5}=0{,}8=80\%$.\par
\elegxos $0{,}2+0{,}8=1$, και $1-0{,}2=0{,}8$.

\paradeigma{Η τράπουλα}
\tixero Μια τράπουλα έχει $52$ φύλλα: τέσσερις σειρές (κούπες, καρό, σπαθιά, μπαστούνια) των $13$ φύλλων, με άσο, $2$ ως $10$, βαλέ, ντάμα και ρήγα. Τραβάμε ένα φύλλο τυχαία.\par
\tithelo Την πιθανότητα: α) κούπα, β) άσος, γ) φιγούρα (βαλές, ντάμα ή ρήγας).\par
\possindeo Υπάρχουν $52$ ισοπίθανα αποτελέσματα.\par
\apantisi α) $\dfrac{13}{52}=\dfrac{1}{4}=25\%$. β) $\dfrac{4}{52}=\dfrac{1}{13}\approx0{,}077\approx7{,}7\%$. γ) $3\cdot4=12$ φιγούρες, άρα $\dfrac{12}{52}=\dfrac{3}{13}\approx0{,}231\approx23{,}1\%$.\par
\elegxos Στο α) ένα στα τέσσερα φύλλα είναι κούπα, όσες οι σειρές.

\paradeigma{Δύο ζάρια}
\tixero Ρίχνουμε δύο ζάρια και προσθέτουμε τα αποτελέσματα.\par
\tithelo Την πιθανότητα το άθροισμα να είναι $7$, να είναι $12$ και να έχουμε το ίδιο αποτέλεσμα στα δύο ζάρια.\par
\possindeo Ισοπίθανα είναι τα $36$ ζευγάρια, όχι τα $11$ αθροίσματα.\par
\apantisi Άθροισμα $7$: $(1,6),(2,5),(3,4),(4,3),(5,2),(6,1)$, $6$ ζευγάρια, άρα $\dfrac{6}{36}=\dfrac{1}{6}\approx16{,}7\%$. Άθροισμα $12$: μόνο το $(6,6)$, άρα $\dfrac{1}{36}\approx2{,}8\%$. Ίδιο αποτέλεσμα: $(1,1),\dots,(6,6)$, $6$ ζευγάρια, άρα $\dfrac{6}{36}=\dfrac{1}{6}$.\par
\elegxos Το άθροισμα $7$ είναι έξι φορές πιο πιθανό από το $12$, όπως δείχνει ο πίνακας της προηγούμενης παραγράφου.

\paradeigma{Ο τροχός}
\tixero Ο τροχός του σχήματος έχει $8$ ίσα τμήματα: $3$ κόκκινα, $4$ μπλε, $1$ κίτρινο.\par
\tithelo Την πιθανότητα να σταματήσει σε κόκκινο, μπλε, κίτρινο.\par
\possindeo Ισοπίθανα είναι τα $8$ τμήματα.\par
\apantisi $\mathrm{Ρ}(\text{κόκκινο})=\dfrac{3}{8}=0{,}375=37{,}5\%$. $\mathrm{Ρ}(\text{μπλε})=\dfrac{4}{8}=\dfrac{1}{2}=50\%$. $\mathrm{Ρ}(\text{κίτρινο})=\dfrac{1}{8}=0{,}125=12{,}5\%$.\par
\elegxos $\dfrac{3}{8}+\dfrac{4}{8}+\dfrac{1}{8}=\dfrac{8}{8}=1$. Κάπου σταματά σίγουρα, άρα το άθροισμα είναι $1$.

\paradeigma{Σχετική συχνότητα}
\tixero Ο Ηλίας έριξε ένα ζάρι $60$ φορές και έφερε $6$ δώδεκα φορές.\par
\tithelo Τη σχετική συχνότητα του $6$ και να τη συγκρίνουμε με την πιθανότητα.\par
\possindeo Σχετική συχνότητα $=\kappa:\nu$.\par
\apantisi $\dfrac{12}{60}=\dfrac{1}{5}=0{,}2=20\%$. Η πιθανότητα είναι $\dfrac{1}{6}\approx16{,}7\%$. Δεν είναι ίσες, αλλά είναι κοντά.\par
\elegxos Με $60$ ρίψεις δεν περιμένουμε ακριβώς $\frac{1}{6}$. Αν ο Ηλίας έριχνε το ζάρι $6000$ φορές, θα περιμέναμε σχετική συχνότητα πολύ πιο κοντά στο $0{,}167$.

\tmima{Ασκήσεις}

\zoni{Το κατάλαβα}{ενδεχόμενα και κλάσματα}

\ex Ρίχνουμε ένα ζάρι. Βρες την πιθανότητα: α) να φέρω $5$, β) να φέρω άρτιο, γ) να φέρω αριθμό μικρότερο από $3$, δ) να μη φέρω $6$.

\ex Ένα κουτί έχει $8$ μπάλες: $4$ κόκκινες, $3$ μπλε και $1$ κίτρινη. Βγάζουμε μία τυχαία. Βρες ως κλάσμα, δεκαδικό και ποσοστό την πιθανότητα να είναι: α) κόκκινη, β) μπλε, γ) όχι κίτρινη, δ) πράσινη.

\ex Συμπλήρωσε τον πίνακα. \par
\begin{center}
\renewcommand{\arraystretch}{1.25}
\begin{tabular}{lcccccc}
\toprule
Κλάσμα & $\frac{1}{2}$ & $\ldots$ & $\ldots$ & $\frac{3}{5}$ & $\ldots$ & $\ldots$\\
Δεκαδικός & $\ldots$ & $0{,}25$ & $\ldots$ & $\ldots$ & $0{,}05$ & $\ldots$\\
Ποσοστό & $\ldots$ & $\ldots$ & $30\%$ & $\ldots$ & $\ldots$ & $100\%$\\
\bottomrule
\end{tabular}
\end{center}

\ex Σωστό ή λάθος; α) Η πιθανότητα ενός ενδεχομένου μπορεί να είναι $1{,}5$. β) Το αδύνατο ενδεχόμενο έχει πιθανότητα $0$. γ) Αν $\mathrm{Ρ}(\mathrm{Α})=0{,}3$, τότε η πιθανότητα να μη συμβεί το $\mathrm{Α}$ είναι $0{,}7$. δ) Το βέβαιο ενδεχόμενο έχει πιθανότητα $100\%$. ε) Στη ρίψη ενός ζαριού η πιθανότητα να φέρω $7$ είναι $\frac{1}{6}$. στ) Αν η πιθανότητα ενός ενδεχομένου είναι $\frac{1}{2}$, τότε σε $10$ επαναλήψεις θα συμβεί ακριβώς $5$ φορές.

\ex Πες αν το ενδεχόμενο είναι βέβαιο, αδύνατο ή ούτε το ένα ούτε το άλλο. α) Ρίχνω ζάρι και φέρνω αριθμό μικρότερο από $7$. β) Ρίχνω ζάρι και φέρνω $0$. γ) Ρίχνω νόμισμα και φέρνω Κεφάλι. δ) Τραβάω ένα φύλλο από τράπουλα (χωρίς τζόκερ) και είναι «$15$ καρό». ε) Από κουτί που έχει μόνο κόκκινες μπάλες, βγάζω κόκκινη.

\ex Επιλέγουμε τυχαία ένα γράμμα της λέξης ΜΑΘΗΜΑΤΙΚΑ (κάθε θέση της λέξης μετράει). Βρες την πιθανότητα να είναι: α) το Α, β) το Μ, γ) φωνήεν, δ) το Ω.

\ex Ρίχνουμε δύο ζάρια. Βρες την πιθανότητα: α) το άθροισμα να είναι $8$, β) να φέρουμε το ίδιο αποτέλεσμα, γ) το άθροισμα να είναι $1$.

\zoni{Το χρησιμοποιώ}{κληρώσεις, τροχοί, τράπουλα}

\ex \textbf{Η κλήρωση.} Στη γιορτή του σχολείου πουλήθηκαν $200$ λαχνοί και κερδίζουν οι $10$. Ποια είναι η πιθανότητα να κερδίσει ένας λαχνός και ποια να μην κερδίσει; Γράψε τες ως κλάσμα, δεκαδικό και ποσοστό.

\ex \textbf{Η τάξη.} Στην Α3 είναι $25$ μαθητές, $14$ κορίτσια και $11$ αγόρια. Επιλέγουμε τυχαία έναν μαθητή. Βρες την πιθανότητα να είναι κορίτσι και να είναι αγόρι.

\ex \textbf{Ο τροχός.} Ένας τροχός έχει $12$ ίσα τμήματα: $5$ κόκκινα, $4$ μπλε, $2$ πράσινα και $1$ χρυσό. Βρες την πιθανότητα για κάθε χρώμα. Πόσο είναι το άθροισμά τους και γιατί;

\ex \textbf{Ο κωδικός.} Ο Νίκος ξέχασε το τελευταίο ψηφίο του κωδικού του κινητού του (ψηφίο από $0$ ως $9$), αλλά ξέρει ότι είναι περιττός. Αν μαντέψει ένα περιττό ψηφίο, ποια είναι η πιθανότητα να το βρει;

\ex \textbf{Η τράπουλα.} Τραβάμε ένα φύλλο από μια τράπουλα $52$ φύλλων. Βρες την πιθανότητα: α) να είναι κόκκινος άσος, β) να είναι φιγούρα, γ) να μην είναι φιγούρα.

\ex \textbf{Οι μπάλες.} Ένα κουτί έχει $20$ μπάλες κόκκινες, μπλε και πράσινες. Η πιθανότητα να βγει κόκκινη είναι $0{,}35$ και να βγει μπλε $\frac{2}{5}$. Πόσες είναι οι μπάλες κάθε χρώματος;

\ex \textbf{Τρεις ρίψεις.} Ρίχνουμε ένα νόμισμα τρεις φορές. Με δενδροδιάγραμμα βρες την πιθανότητα: α) να φέρουμε Κεφάλι και τις τρεις φορές, β) να φέρουμε Γράμματα τουλάχιστον μία φορά, γ) να φέρουμε Κεφάλι ακριβώς δύο φορές.

\zoni{Το σκέφτομαι}{τι σημαίνει «πιθανό»}

\ex Ο Στέλιος λέει: «Αύριο ή θα βρέξει ή δεν θα βρέξει, άρα η πιθανότητα να βρέξει είναι $50\%$». Συμφωνείς; Σύγκρινε με τη ρίψη νομίσματος.

\ex Ο Ηλίας έριξε ένα ζάρι $20$, $100$ και $600$ φορές και έφερε $6$ αντίστοιχα $5$, $18$ και $98$ φορές. Βρες τις σχετικές συχνότητες του $6$ και σύγκρινέ τες με την πιθανότητα $\frac{1}{6}\approx0{,}167$. Τι παρατηρείς όσο αυξάνονται οι ρίψεις;

\ex Σχεδίασε έναν τροχό με $8$ ίσα τμήματα, ώστε η πιθανότητα για κόκκινο να είναι $\frac{3}{8}$ και για μπλε $\frac{1}{4}$. Πόσα τμήματα έμειναν για άλλα χρώματα;

\ex Ο Νίκος λέει: «Στα δύο ζάρια τα αθροίσματα είναι $11$ ($2$ ως $12$), άρα η πιθανότητα για άθροισμα $7$ είναι $\frac{1}{11}$». Βρες το λάθος και τη σωστή πιθανότητα.

\zoni{Για το σπίτι}{πείραμα και πρόβλημα}

\ex Ρίξε ένα νόμισμα $20$ φορές και γράψε πόσες φορές ήρθε Κεφάλι. Βρες τη σχετική συχνότητα και σύγκρινέ τη με την πιθανότητα $\frac{1}{2}$. Αν μπορείς, σύγκρινε τα αποτελέσματά σου με αυτά ενός φίλου.

\ex Φτιάξε ένα δικό σου πρόβλημα πιθανότητας με κουτί και χρωματιστές μπάλες, με τράπουλα ή με ζάρι. Λύσε το και γράψε την πιθανότητα ως κλάσμα, δεκαδικό και ποσοστό.
